% Shared scientific content. Keep scientific wording in this file.
\newcommand{\SharedScientificBody}{%
\section{Introduction}
\label{sec:introduction}

The granular approach restored an elementary but consequential fact to
macroeconomic aggregation: units enter an aggregate according to their
economic weights, so a small number of very large firms can prevent
idiosyncratic fluctuations from averaging out rapidly
\citep{gabaix2011}. Gabaix's scaling result combines two logically
distinct steps. First, assumptions on the firm-size distribution
determine the population path of concentration,
\begin{equation}
\label{eq:gabaix_concentration}
    H(N)=\sum_{i=1}^{N}w_i(N)^2,
    \qquad
    h(N)=H(N)^{1/2}.
\end{equation}
Second, homogeneous marginal shock volatility and mutually uncorrelated
firm-level shocks map that concentration path into aggregate volatility,
\begin{align}
\label{eq:gabaix_volatility_bridge}
    \sigma_X(N) &= \sigma h(N), \\
    V(N) &= \sigma^2H(N).
\end{align}
The resulting population-scaling law is exact within that benchmark.
The firm-size distribution determines the concentration side of the
result; the homogeneous and uncorrelated shock structure supplies the
stochastic bridge from concentration to aggregate volatility.

We work primarily with aggregate variance because its diagonal and
covariance components are additive. Since
$\sigma_X(N)=V(N)^{1/2}$, the standard-deviation and variance
elasticities satisfy
\begin{equation}
\label{eq:variance_volatility_bridge}
    -\frac{d\log \sigma_X(N)}{d\log N}
    =
    \frac{1}{2}
    \left[
        -\frac{d\log V(N)}{d\log N}
    \right].
\end{equation}
Likewise, the elasticity of $h(N)$ is one half of the elasticity of
$H(N)$. The distinction developed below is therefore invariant to the
choice of variance or standard-deviation units, provided this factor of
two is made explicit.

Concentration and scaling are nevertheless different objects.
Concentration records the distribution of direct economic weights at a
given scale. Aggregate-variance scaling records how a complete
population-indexed quadratic form changes along a specified path. For
the weighted aggregate
\begin{equation}
    X_t(N)=\sum_{i=1}^{N}w_i(N)\varepsilon_{i,t}(N),
\end{equation}
aggregate variance is
\begin{align}
\label{eq:intro_variance}
    V(N)
    &=
    \sum_{i=1}^{N}w_i(N)^2
    \operatorname{Var}\!\left[\varepsilon_{i,t}(N)\right]
    \nonumber\\
    &\quad+
    2\sum_{i<j}w_i(N)w_j(N)
    \operatorname{Cov}\!\left[
        \varepsilon_{i,t}(N),\varepsilon_{j,t}(N)
    \right].
\end{align}
Writing $\mathbf{w}(N)$ for the direct-weight vector and
$\Sigma_{\varepsilon}(N)$ for the covariance matrix of the observed
contributions,
\begin{equation}
\label{eq:matrix_variance}
    V(N)
    =
    \mathbf{w}(N)^{\mathsf T}
    \Sigma_{\varepsilon}(N)
    \mathbf{w}(N),
    \qquad
    H(N)=\mathbf{w}(N)^{\mathsf T}\mathbf{w}(N).
\end{equation}
The granular benchmark obtains equation~\eqref{eq:gabaix_volatility_bridge}
from the sufficient closure
$\Sigma_{\varepsilon}(N)=\sigma^2 I$. Even under independence, allowing
the common marginal variance to vary with population scale gives
\begin{equation}
\label{eq:benchmark_elasticity_bridge}
    -\frac{d\log V(N)}{d\log N}
    =
    -\frac{d\log H(N)}{d\log N}
    -\frac{d\log \sigma^2(N)}{d\log N}.
\end{equation}
The concentration elasticity then identifies the variance elasticity
only when the marginal-variance term is scale neutral. Representative
firm-level evidence instead finds that growth volatility varies
systematically with firm size, and therefore cannot generally be
compressed into one common, scale-invariant scalar
\citep{calvinoEtAl2018}.

The scalar-identity closure is sufficient but not necessary. In
multisector models, independent primitive innovations may be propagated
into correlated sectoral outcomes. If $\mathbf{y}=G\boldsymbol{\eta}$
and $\Omega=\operatorname{Var}(\boldsymbol{\eta})$, then
\begin{equation}
\label{eq:effective_exposure_representation}
    V
    =
    \mathbf{w}^{\mathsf T}G\Omega G^{\mathsf T}\mathbf{w}
    =
    \mathbf{a}^{\mathsf T}\Omega\mathbf{a},
    \qquad
    \mathbf{a}=G^{\mathsf T}\mathbf{w}.
\end{equation}
Economic linkages can therefore be represented either as induced
covariance among observed components or as transformed effective
exposures to primitive shocks. \citet{dupor1999} gives conditions under
which a nontrivial propagation structure preserves benchmark
second-moment scaling along the aggregate direction, whereas
\citet{horvath1998} shows that a different population path for supplier
structure can slow diversification even with equal direct weights and
independent primitive shocks. The relevant restriction is therefore not
the absence of every off-diagonal term, but the existence of a
scale-neutral bridge between direct weights and aggregate stochastic
exposure.

A related firm-growth literature studies a different estimand (the
volatility of a firm conditional on its size) but reaches a parallel
accounting lesson. Nonclassical exponents can arise from heterogeneous
independent constituent sizes \citep{sutton2001}, from the sublinear
expansion of diversification breadth \citep{bottazziSecchi2006}, or from
a long crossover induced by heterogeneous product counts and sizes
\citep{riccaboniEtAl2008}. \citet{calvinoEtAl2018} further show that the
same reduced-form firm size--volatility exponent can be generated either
by scale-dependent correlation among subunits or by a sublinear relation
between firm size and subunit count. These results establish that a
fitted exponent is conditional on constituent and dependence structure;
they do not make the exponent a unique diagnostic of that structure.

This paper does not dispute the conditional algebra of the granular
benchmark or these model-specific scaling mechanisms. It asks a narrower
identification question: what follows from the complete direct-weight
path when the scale behavior of marginal variances, effective exposures,
and dependence has not already been fixed? We answer in two steps.

First, we show that identical complete concentration paths can generate
different aggregate-variance scaling over admissible covariance paths.
Under the equal-weight IID benchmark, $H(N)=V(N)=1/N$. Holding the same
weights and marginal variances fixed while changing only dependence can
instead produce a positive variance floor and an elasticity converging
to zero. The equicorrelation construction is a logical witness for
non-identification, not a proposed empirical model of firms.

Second, we derive an exact finite-interval attribution of aggregate
variance elasticity to changes in diagonal variance and signed
covariance components. The formula never takes logarithms of individual
components and therefore remains well defined when covariance
contributions are negative, zero, or sign-changing. In the
differentiable positive-component case, the result reduces to a local
contribution-share-weighted identity.

The mathematical ingredients are established. Standard variance
accounting supplies diagonal and covariance components, while
logarithmic-mean methods supply residual-free exact-change accounting
\citep{ang2005}. Extensions of logarithmic-mean factor decomposition also
address zero, negative, and sign-changing values through special
procedures \citep{angLiu2007}. Our claim is correspondingly narrow: we
specialize the logarithmic-mean bridge to the finite-interval elasticity
of population-indexed aggregate variance, treating signed covariance as
an additive component whose level is never logged. This specialization
separates a level statement about direct concentration from a path
statement about stochastic exposure and yields an accounting identity
that can be checked directly in empirical or synthetic data.

\section{Exact decomposition}
\label{sec:exact_decomposition}

Consider an aggregate composed of $Q$ contribution series,
\begin{equation}
    X_t(N)=\sum_{q=1}^{Q}c_{q,t}(N),
\end{equation}
where $N>0$ indexes population scale. Its variance is
\begin{equation}
\label{eq:variance_accounting}
    V(N)
    =
    \sum_{q=1}^{Q}\operatorname{Var}\!\left[c_{q,t}(N)\right]
    +
    2\sum_{q<r}\operatorname{Cov}\!\left[c_{q,t}(N),c_{r,t}(N)\right].
\end{equation}
Let $\mathcal{K}$ collect the diagonal and doubled
off-diagonal terms,
\begin{align}
    C_{qq}(N)
    &=
    \operatorname{Var}\!\left[c_{q,t}(N)\right], \\
    C_{qr}(N)
    &=
    2\operatorname{Cov}\!\left[c_{q,t}(N),c_{r,t}(N)\right],
    \qquad q<r.
\end{align}
Then
\begin{equation}
\label{eq:component_sum}
    V(N)=\sum_{k\in\mathcal{K}}C_k(N).
\end{equation}
Diagonal components are non-negative, whereas covariance components may
have either sign. The component universe $\mathcal{K}$ must retain a
comparable meaning at the population levels being contrasted. When
units enter or exit, both endpoints must be embedded in a common index,
for example by zero-filling absent contribution terms or by aggregating
them into stable, consistently defined groups. Otherwise the attribution
mixes within-component scaling with entry, exit, or reclassification.

For $N_0,N_1>0$, $N_0\neq N_1$, and positive aggregate variance at both
levels, define the finite-interval variance elasticity
\begin{equation}
\label{eq:direct_alpha}
    \widehat{\alpha}(N_0,N_1)
    =
    -\frac{
        \log V(N_1)-\log V(N_0)
    }{
        \log N_1-\log N_0
    }.
\end{equation}
For $a,b>0$, let
\begin{equation}
\label{eq:log_mean}
    L(a,b)
    =
    \begin{cases}
        \displaystyle\frac{a-b}{\log a-\log b}, & a\neq b,\\[1em]
        a, & a=b,
    \end{cases}
\end{equation}
be the logarithmic mean. This is the same exact-change bridge used by
logarithmic-mean Divisia decomposition methods \citep{ang2005}; the
object specialized below is a population elasticity of an additive
variance--covariance accounting system.

\begin{proposition}[Exact finite-interval decomposition]
\label{prop:exact_decomposition}
For each $k\in\mathcal{K}$, define
\begin{equation}
\label{eq:component_alpha_contribution}
    \widehat{\alpha}^{\,C}_k(N_0,N_1)
    =
    -\frac{
        C_k(N_1)-C_k(N_0)
    }{
        L\!\left(V(N_1),V(N_0)\right)
        \left[\log N_1-\log N_0\right]
    }.
\end{equation}
Then
\begin{equation}
\label{eq:exact_alpha_decomposition}
    \widehat{\alpha}(N_0,N_1)
    =
    \sum_{k\in\mathcal{K}}
    \widehat{\alpha}^{\,C}_k(N_0,N_1).
\end{equation}
The identity remains valid when individual covariance components are
negative.
\end{proposition}

\begin{proof}
Equation~\eqref{eq:component_sum} implies
\begin{equation}
    \sum_{k\in\mathcal{K}}
    \left[C_k(N_1)-C_k(N_0)\right]
    =V(N_1)-V(N_0).
\end{equation}
Summing equation~\eqref{eq:component_alpha_contribution} over $k$ and
using
\begin{equation}
    \frac{V(N_1)-V(N_0)}{L(V(N_1),V(N_0))}
    =\log V(N_1)-\log V(N_0)
\end{equation}
produces equation~\eqref{eq:direct_alpha} exactly.
\end{proof}

The proposition is a specialization of logarithmic-mean exact-change
accounting rather than a new general decomposition calculus. Its role is
to translate the additive variance identity into the finite-interval
population elasticity of interest. Only total variance is logged.
Individual covariance components may therefore be negative, zero, or
sign-changing, provided aggregate variance is positive at both
endpoints. This differs from multiplicative factor decompositions, where
negative or sign-changing factors require additional conventions or
analytical-limit procedures \citep{angLiu2007}.

The result is exact for each finite interval and does not require a
global power law. A regression coefficient estimated over many
population levels may still provide a useful summary, but it averages
interval-specific elasticities whose component contributions can change
across the fitted range. This distinction matters in settings where an
apparently stable fitted exponent may conceal a slow crossover
\citep{riccaboniEtAl2008}. The decomposition reconstructs the accounting
sources of the observed change; it does not identify the causal mechanism
that generated those component paths.

The attribution is also consistent under grouping. For any group
$G\subseteq\mathcal{K}$, define
$C_G(N)=\sum_{k\in G}C_k(N)$ and assign it the contribution
\begin{equation}
\label{eq:grouped_alpha_contribution}
    \widehat{\alpha}^{\,C}_G(N_0,N_1)
    =
    -\frac{C_G(N_1)-C_G(N_0)}{
        L\!\left(V(N_1),V(N_0)\right)
        [\log N_1-\log N_0]
    }.
\end{equation}
Then
$\widehat{\alpha}^{\,C}_G=\sum_{k\in G}\widehat{\alpha}^{\,C}_k$.
Pairwise covariance terms may therefore be aggregated into stable
within-sector, between-sector, or other substantively defined groups
without changing the reconstructed total.

In the differentiable limit,
\begin{equation}
\label{eq:local_signed_decomposition}
    \alpha_{\mathrm{eff}}(N)
    =-
    \frac{d\log V(N)}{d\log N}
    =-
    \sum_{k\in\mathcal{K}}
    \frac{N C_k'(N)}{V(N)}.
\end{equation}
If every component is positive, define
\begin{equation}
    \alpha_k(N)
    =
    -\frac{d\log C_k(N)}{d\log N}.
\end{equation}
Then
\begin{equation}
\label{eq:weighted_local_alpha}
    \alpha_{\mathrm{eff}}(N)
    =
    \sum_{k\in\mathcal{K}}
    \frac{C_k(N)}{V(N)}\alpha_k(N).
\end{equation}
At each $N$, equation~\eqref{eq:weighted_local_alpha} is a
contribution-share-weighted identity for local elasticities; it is not a
fixed finite-interval average. Its weights generally vary with $N$. A
fixed average arises only under additional restrictions, such as common
component elasticities or invariant contribution shares. When positive
components decay at different power-law rates, their shares change with
$N$ and the slowest-decaying component eventually dominates. This broad
component-balance intuition is already present in constituent-growth and
scale-dependent-correlation models \citep{sutton2001,riccaboniEtAl2008,calvinoEtAl2018}; the result here supplies an exact signed accounting
representation for a different population-indexed estimand.

With signed covariance components, normalized contributions need not be
non-negative, so the analogous representation is not a convex
combination and the aggregate elasticity need not lie within the range
of the diagonal elasticities. Covariance is therefore not merely a
correction to the level of aggregate variance; its own scale path can
reinforce, offset, or leave unchanged the diversification generated by
diagonal terms.

Comparability across scales is substantive. If parts appear, disappear,
or change meaning between $N_0$ and $N_1$, the measured elasticity may
combine diversification with composition, partition, or dependence
changes. The algebra remains exact only after those changes are
represented through a common accounting structure; the economic
interpretation then depends on whether entry, exit, and reclassification
are intended components of the population path. Cross-sectional
size--volatility relations such as those studied by
\citet{sutton2001,bottazziSecchi2006,riccaboniEtAl2008,calvinoEtAl2018}
are therefore neighboring scaling objects, not automatically
component-matched population paths of the type required here.

\section{Why concentration does not identify variance scaling}
\label{sec:concentration_scaling}

Consider the weighted aggregate
\begin{equation}
\label{eq:weighted_aggregate}
    X_t(N)
    =
    \sum_{i=1}^{N}w_i(N)\varepsilon_{i,t}(N),
    \qquad
    \sum_{i=1}^{N}w_i(N)=1.
\end{equation}
For equal weights, $w_i(N)=1/N$, the concentration path is
\begin{equation}
\label{eq:equal_weight_concentration}
    H(N)=\frac{1}{N}.
\end{equation}
Suppose first that shocks have unit variance and are mutually
independent. Aggregate variance is then
\begin{equation}
\label{eq:iid_variance}
    V_{\mathrm{iid}}(N)
    =
    \operatorname{Var}\!\left[
        \frac{1}{N}\sum_{i=1}^{N}\varepsilon_{i,t}
    \right]
    =\frac{1}{N},
\end{equation}
and its local variance elasticity is
\begin{equation}
    \alpha_{\mathrm{iid}}(N)=1.
\end{equation}

Now retain the same weights and marginal variances but impose a common
pairwise correlation $\rho\in(0,1)$. Aggregate variance becomes
\begin{align}
\label{eq:correlated_variance}
    V_{\rho}(N)
    &=
    \frac{1}{N^2}\left[N+\rho N(N-1)\right]
    \nonumber\\
    &=
    \rho+\frac{1-\rho}{N},
\end{align}
with local elasticity
\begin{equation}
\label{eq:correlated_alpha}
    \alpha_{\rho}(N)
    =
    \frac{1-\rho}{\rho N+1-\rho}.
\end{equation}
The independent aggregate continues to diversify at rate $N^{-1}$,
whereas the correlated aggregate approaches the variance floor
$\lim_{N\to\infty}V_{\rho}(N)=\rho$ and
$\lim_{N\to\infty}\alpha_{\rho}(N)=0$. Neither covariance path is offered
as a complete empirical model. Their purpose is to isolate the missing
information by holding the entire direct-weight path and all marginal
variances fixed.

\begin{figure*}[t]
    \centering
    \includegraphics[
        width=0.98\textwidth
    ]{figures/figure_1_concentration_is_not_scaling.png}
    \caption{
    Concentration and aggregate-variance scaling under equal weights
    and unit marginal shock variance. Panel~(a) compares the common
    concentration path $H(N)=1/N$ with aggregate variance under
    independence and under equicorrelation $\rho=0.15$. The independent
    variance coincides with $H(N)$, whereas the correlated variance
    approaches the floor $\rho$. Panel~(b) reports adjacent-interval
    variance elasticities: concentration and independent variance retain
    elasticity one, while the correlated variance elasticity falls
    toward zero.
    }
    \label{fig:concentration_scaling}
\end{figure*}

The counterexample establishes non-identification, not merely a
difference in variance levels. On the displayed grid, the correlated
adjacent-interval variance elasticity falls from approximately $0.500$
to $0.113$, while the concentration elasticity remains one. The two
systems have the same population, weights, Herfindahl index, effective
number of firms, and complete concentration path. What differs is the
population-indexed variance--covariance structure.

Related multisector results provide structural antecedents to the same
broad warning. Under an aggregate-direction symmetry,
\citet{dupor1999} obtains benchmark diversification despite nontrivial
propagation, whereas \citet{horvath1998} obtains slower scaling when the
number of broadly supplying sectors grows sublinearly with total sector
count. These papers do not prove the fixed-weight covariance
counterexample or the finite-interval decomposition developed here, but
they show that direct sector count and direct weights do not by
themselves determine how stochastic exposure scales.

A fitted concentration law can therefore identify aggregate-variance
scaling only after marginal variance, effective exposure, dependence,
and composition have been shown to provide a scale-neutral bridge. This
qualifies rather than overturns the granular mechanism: unequal weights
may remain quantitatively important, but the firm-size distribution is
not by itself sufficient to identify the population elasticity of
aggregate variance.

\section{Diversification as a component balance}
\label{sec:component_balance}

The independent homogeneous benchmark is a diagonal-dilution result.
For an equal-weight average of $N$ mutually independent unit-variance
contributions, aggregate variance falls as $1/N$. More generally, the
observed elasticity records how every source of variance changes along
the population path. The surrounding literature already supplies
several mechanism-specific examples: heterogeneous independent
constituents \citep{sutton2001,riccaboniEtAl2008}, diversification breadth
\citep{bottazziSecchi2006}, and scale-dependent correlation
\citep{calvinoEtAl2018}. A value below one therefore need not represent
one common mechanism; it may arise because diagonal components decay at
different rates, because covariance reinforces or offsets their decline,
or because the composition of the aggregate changes with $N$.

The first attribution example uses two independent positive components,
\begin{equation}
\label{eq:heterogeneous_demo}
    V_{\mathrm{het}}(N)
    =
    N^{-0.9}+0.4N^{-0.3}.
\end{equation}
This example generalizes the common-exponent inheritance case. If all positive components share the same elasticity, the
aggregate inherits that elasticity. When
component elasticities differ, however, no fixed population exponent
is primitive. The aggregate elasticity changes as the components'
variance shares change. In Figure~\ref{fig:alpha_verification}(a), the
fast-decaying component initially supplies most of the variance decline,
but the slow component becomes increasingly important; the aggregate
elasticity consequently falls from approximately $0.582$ to $0.422$.

The second attribution example groups the equicorrelated counterexample
into total diagonal variance and total doubled covariance,
\begin{align}
\label{eq:grouped_correlated_demo}
    D(N) &= \frac{1}{N}, \\
    B(N) &= \rho\left(1-\frac{1}{N}\right), \\
    V_{\rho}(N) &= D(N)+B(N).
\end{align}
Although $B(N)$ is positive, it rises with population size. Hence
$B(N_1)-B(N_0)>0$ for $N_1>N_0$, and the minus sign in
equation~\eqref{eq:component_alpha_contribution} makes its contribution
to the variance-decay elasticity negative. This sign is path dependent,
not a consequence of positive covariance alone. A declining positive
covariance component would contribute positively to variance decay, and
a scale-neutral dependence structure may leave the benchmark rate
unchanged, as the contrast between \citet{dupor1999} and
\citet{calvinoEtAl2018} helps illustrate.

\begin{figure*}[t]
    \centering
    \includegraphics[
        width=0.98\textwidth
    ]{figures/figure_2_exact_component_attribution.png}
    \caption{
    Exact component attribution of finite-interval variance
    elasticities. Panel~(a) decomposes the heterogeneous independent
    aggregate into fast- and slow-decaying diagonal contributions.
    Panel~(b) decomposes the correlated equal-weight aggregate into
    diagonal and total doubled-covariance contributions. The positive
    covariance level increases along the population path and therefore
    contributes negatively to variance decay. Open reconstructed-total
    markers overlap the directly computed totals up to floating-point
    precision. The numerical annotation reports the maximum residual
    over the two displayed panels rather than over the complete
    verification artifact.
    }
    \label{fig:alpha_verification}
\end{figure*}

In this equal-weight case, the local decomposition has the especially
transparent form
\begin{equation}
\label{eq:diversification_balance}
    \alpha_V(N)
    =
    \frac{D(N)}{V(N)}
    -
    \frac{N B'(N)}{V(N)}.
\end{equation}
The first term records diagonal variance dilution; the second records
the signed effect of the changing covariance level. Independence gives
$B(N)=0$ and restores $\alpha_V(N)=1$. Under equicorrelation,
$B'(N)>0$, so covariance accumulation offsets diagonal diversification
and pushes the aggregate elasticity toward zero.

Equation~\eqref{eq:diversification_balance} gives a precise meaning to
delayed diversification. The scaling exponent is not itself a primitive
mechanism, nor is it a synonym for concentration. It is the net
accounting outcome of diagonal dilution and covariance accumulation
along the population path. Figure~\ref{fig:alpha_verification}(b)
displays this balance: the diagonal contribution is positive, while the
increasing positive covariance level makes a negative contribution and can offset variance decay.

The accompanying machine-readable artifact evaluates these identities
on $N\in\{4,8,16,32,64\}$. At each population level, direct and
reconstructed variance are generated from the same component sum, so
their agreement verifies table construction rather than providing an
independent mathematical test. The adjacent-interval checks separately
compare the direct log-ratio elasticity with the sum of signed component
contributions. All pass the declared scale-aware tolerance contract, and
the maximum absolute scaling-reconstruction residual over the complete
reference artifact is $6.7\times10^{-16}$. Last-digit residuals may vary
across numerical environments; the scientific criterion is the declared
tolerance. The exercise verifies implementation and internal consistency,
not the scaling behavior of any particular economy.

\section{Conclusion}
\label{sec:conclusion}

A concentration path describes how direct economic weight is
distributed. Aggregate-variance scaling describes how a complete
population-indexed stochastic exposure,
$\mathbf{w}(N)^{\mathsf T}\Sigma_{\varepsilon}(N)\mathbf{w}(N)$ or its
effective-exposure equivalent, changes with population scale. The two
coincide under the homogeneous, mutually uncorrelated, scale-invariant
benchmark of \citet{gabaix2011}, and more generally whenever the
stochastic structure supplies a scale-neutral bridge along the relevant
aggregate direction \citep{dupor1999}. Without such restrictions, even
the complete direct-weight path does not identify aggregate variance or
its elasticity; network structure, constituent organization, marginal
variance, and dependence may each alter the result
\citep{horvath1998,sutton2001,bottazziSecchi2006,riccaboniEtAl2008,calvinoEtAl2018}.

The surrounding literature therefore contains many valid conditional
scaling mechanisms. The present result does not replace them or infer a
unique mechanism from an observed exponent. It specializes inherited
logarithmic-mean exact-change accounting \citep{ang2005,angLiu2007} to a
different task: exact finite-interval attribution of aggregate variance
elasticity to comparable diagonal and signed covariance changes. An
observed exponent is thereby represented as the net accounting outcome
of component paths. Common elasticities reproduce familiar inheritance
results; heterogeneous elasticities shift contribution shares; and
covariance may reinforce, offset, or leave unchanged diagonal
diversification depending on its own scale path.

This distinction changes the order of empirical analysis. A fitted
concentration or volatility exponent may be consistent with several
latent structures, so reproducing it is evidence of possibility rather
than unique structural identification. Prior disaggregated work already
illustrates a diagonal component-accounting strategy at the firm and
submarket level \citep{bottazziSecchi2006}. The framework here extends
that accounting contract to signed covariance: define a coherent
population path, construct comparable components at its endpoints,
estimate the aggregate variance elasticity directly, and verify that the
component contributions reconstruct it before assigning a causal
mechanism. When internal firm structure is unobserved, this procedure can
still be applied at the observed firm level, but it cannot by itself
explain the marginal variance generated inside each firm.


% Declarations


\section*{Funding}

This research did not receive any specific grant from funding agencies
in the public, commercial, or not-for-profit sectors.

\section*{Declaration of competing interest}

The author declares that he has no known competing financial interests
or personal relationships that could have appeared to influence the
work reported in this paper.

\section*{Data and code availability}

No empirical data were used in this study. Source code, deterministic
computational constructions, tests, and generated verification artifacts
are available at the following \href{https://github.com/matuteiglesias/concentration-is-not-scaling}{Github repository} 

\section*{
Declaration of generative AI and AI-assisted technologies in the
writing process
}

During the preparation of this work, the author used OpenAI ChatGPT and
Codex to assist with research organization, software implementation,
adversarial checking, and language revision. The author reviewed and
verified the mathematical derivations, computational outputs,
references, and final text, edited the resulting material as needed,
and takes full responsibility for the content of the publication.
}
